\clearpage
\begin{trapbox}
\begin{center}\begin{tikzpicture}[x=.58cm,y=.58cm,font=\small,>=stealth]
\foreach \s in {0,7.6,15.2}{
\begin{scope}[xshift=\s*0.58cm]
\draw[gray,->] (-3.3,0)--(3.5,0);
\draw[gray,->] (0,-2.3)--(0,2.5);
\draw[CStatement,thick] (0,0) ellipse(3 and 2);
\end{scope}}
\draw[teal,very thick] (-2.598076,1)--(2.598076,1);
\fill[orange!80!black] (0,1) circle(2pt) node[above right] {$M$};
\node at (0,3.2) {$x_0=0,\ y_0\ne0$}; \node at (0,-3) {水平：$y=y_0$};
\begin{scope}[xshift=4.408cm]
\draw[teal,very thick] (1.5,-1.732051)--(1.5,1.732051);
\fill[orange!80!black] (1.5,0) circle(2pt) node[above right] {$M$};
\node at (0,3.2) {$y_0=0,\ x_0\ne0$}; \node at (0,-3) {竖直：$x=x_0$};
\end{scope}
\begin{scope}[xshift=8.816cm]
\draw[teal,very thick] (-3,0)--(3,0);
\draw[purple,dashed,thick] (-2.12132,-1.414214)--(2.12132,1.414214);
\fill[orange!80!black] (0,0) circle(2pt) node[above left] {$M=O$};
\node at (0,3.2) {$x_0=y_0=0$}; \node at (0,-3) {中心弦：方向任意};
\end{scope}
\end{tikzpicture}
\end{center}
\begin{description}[style=nextline,leftmargin=2.8em,itemsep=0.5em,labelsep=0.4em]
\item[\textbf{易错 1｜$y_0=0$ 直接套斜率公式。}]
分母为零不能除。若 $x_0\ne0$ 则 $x=x_0$（竖直弦）；若 $x_0=0$ 则 $M=O$，中心弦方向任意。\textbf{最小反例：}椭圆 $x^2/16+y^2/9=1$ 中点 $O$ 可以对应 $y=0$，并非必然竖直。
\item[\textbf{易错 2｜把 $x_0=0$ 误判成竖直弦。}]
当 $x_0=0,y_0\ne0$，$k=0$，弦\textbf{水平}；例如中点 $(0,1)$，直线 $y=1$。
\item[\textbf{易错 3｜长轴换方向后系数颠倒。}]
不要只盯着 $a,b$ 的字母顺序。对 $x^2/A+y^2/B=1$ 永远先写 $x_0\Delta x/A+y_0\Delta y/B=0$，故 $k=-Bx_0/(Ay_0)$；长轴在 $y$ 轴时令 $A=b^2,B=a^2$ 即可。
\item[\textbf{易错 4｜只看点差关系，不检查弦是否真实存在。}]
中点必须满足 $x_0^2/a^2+y_0^2/b^2<1$。\textbf{最小反例：}椭圆 $x^2/16+y^2/12=1$ 中 $M(4,1)$ 在外，代数上能算出斜率 $-3$，但不存在以它为中点的真实弦。边界点也不能是两个不同交点的中点。
\item[\textbf{易错 5｜把内部点 $M$ 当作原椭圆的切点。}]
一般 $M$ 不在原椭圆上，\textbf{原椭圆在 $M$ 处没有切线}。正确说法：若 $M\ne O$，弦线是过 $M$ 的\textbf{同心小椭圆} $x^2/a^2+y^2/b^2=r^2$ 的切线，其中 $r^2=x_0^2/a^2+y_0^2/b^2$。
\item[\textbf{易错 6｜证明中无根据地约去差值。}]
不能在端点不同的条件下直接断定 $x_1-x_2\ne0$：竖直弦的 $\Delta x$ 本来就是 $0$。先保留\textbf{无除法式}，有 $y_0\ne0$ 时再论证 $\Delta x\ne0$，然后才可求斜率。
\end{description}
\medskip
\textbf{\textcolor{CTrap}{考场防错顺序：}}\quad
\fbox{中点在内？}$\ \longrightarrow\ $
\fbox{中心点？}$\ \longrightarrow\ $
\fbox{先写无除法式}$\ \longrightarrow\ $
\fbox{斜率是否存在？}
\end{trapbox}
